\documentclass[a4paper,10pt]{article}

\usepackage[T1]{fontenc}
\usepackage{hyperref}

%opening
\title{Blender JA dotXSI Tools Manual}
\author{BE-Arbiter}

\begin{document}

 \maketitle

 \newpage

 \tableofcontents

 \newpage

 \section{Purpose of this document}

 This manual explains the Blender add-on importing and exporting Softimage dotXSI 3.0 skeletal
 animations (\texttt{.xsi}) as NLA strips, for Jedi Academy animation work. It is meant to be used
 together with the Blender Jedi Academy Plugin Suite (the \texttt{jediacademy} add-on), which reads and
 writes the game's \texttt{.gla} and \texttt{animation.cfg} files. It does not explain how to use
 Blender itself.

 dotXSI 3.0 is the text format written by 3ds Max's dotXSI exporter, which Raven's tools (Carcass)
 compile into \texttt{.gla} animations. Many community animations for Jedi Academy are distributed in
 it, typically made with Ashuradx's CAT rig of the \texttt{\_humanoid} skeleton.

 \section{Changelog}

 \subsection*{1.0.0}
 Initial release:
 \begin{itemize}
  \item dotXSI 3.0 import as NLA strips, onto the existing \texttt{skeleton\_root} or a new armature
  \item Several sequence names separated by dots create stacked strips, one action per name
  \item Export of the selected NLA strips as dotXSI 3.0, concatenated end to end
 \end{itemize}

 \section{Installation}

 Download \texttt{Blender\_JA\_dotXSI\_Tools.zip} from the releases page and install it with the
 ``Install from Disk'' button in Blender's add-on preferences, then enable
 ``JA dotXSI NLA Import/Export (.xsi)''. Blender 4.1 or newer is required.

 The add-on adds two menu entries:
 \begin{itemize}
  \item File $\rightarrow$ Import $\rightarrow$ JA dotXSI NLA Import (.xsi)
  \item File $\rightarrow$ Export $\rightarrow$ JA dotXSI NLA Export (.xsi)
 \end{itemize}

 \section{The dotXSI 3.0 files}

 Only the text flavour of dotXSI 3.x (header \texttt{xsi 0300txt 0032}) is supported. The older 1.x
 files (header \texttt{xsi 0101txt}, e.g. Battlezone II models) are a different, DirectX-like format
 and are rejected with an error, as are binary files.

 Only the skeleton and its animation are read: every \texttt{SI\_Model} becomes a bone, whatever its
 content (\texttt{SI\_Null}, mesh, \ldots). Meshes, materials, cameras and lights are ignored.

 \subsection{Conventions}

 These were verified against Raven's \texttt{\_humanoid.gla}, compiled by Carcass from the same kind
 of files:
 \begin{itemize}
  \item \texttt{SI\_Transform SRT-} values are local to the parent model, \texttt{BASEPOSE-} values are
  in world space. The base pose is used as the rest pose.
  \item Rotations are Euler XYZ in degrees; a local matrix is translation $\times$ rotation $\times$
  scale.
  \item \texttt{SI\_FCurve}s hold the local scaling, rotation and translation per frame. Tangents of
  \texttt{HERMITE}/\texttt{CUBIC} curves are ignored, keys are interpolated linearly.
  \item The files are Y-up, Ghoul 2 and Blender are Z-up: $(x, y, z) \rightarrow (x, -z, y)$.
  \item \texttt{\_humanoid} is scaled by 0.64. With this scale and axis conversion, the file's base
  pose matches the \texttt{.gla}'s base pose, bone axes included (Ghoul 2 bones point along X).
  \item Carcass moves the whole skeleton by the model's origin: every frame of \texttt{\_humanoid.gla}
  has \texttt{model\_root} at $(0, 0, -24)$ --- the player origin is 24 units above the feet --- while
  its base pose stays at 0. The dotXSI files do not contain this offset.
  \item Carcass reads an \texttt{\_always\_} name suffix as a flag and drops it: the file's
  \texttt{face\_always\_} is the \texttt{.gla}'s \texttt{face} bone.
  \item Bone scale is not supported by Ghoul 2 and is ignored. (The CAT rig's face bones are scaled
  by 1.087.)
 \end{itemize}

 \section{Import}

 \subsection{Target armature}

 The animation is applied to the armature named \texttt{skeleton\_root} (the name the
 \texttt{jediacademy} add-on uses), or else to the active armature, matching bones by name. If there
 is none, a \texttt{skeleton\_root} is created from the file's base pose and a warning is shown. The
 import dialog shows which armature will be used.

 The animation is transferred as world-space offsets from the file's base pose, applied on top of
 each bone's own rest pose. So it also works on an armature whose rest orientations differ, such as
 one imported by the \texttt{jediacademy} add-on with its \texttt{\_humanoid} skeleton changes.

 \subsection{Settings}

 \begin{description}
  \item[Bones] Only when a new armature is created: \emph{JKA \_humanoid} creates the 53 bones of
  \texttt{\_humanoid.gla} with its hierarchy (a bone missing from the file is replaced by its nearest
  present ancestor), \emph{All but helpers} every model except effectors, \texttt{skeleton\_root} and
  \texttt{mesh\_root}, \emph{All} every model with the file's hierarchy.
  \item[Scale] Scale applied to the file, 0.64 for \texttt{\_humanoid}.
  \item[Origin Offset] Translation added to the whole skeleton in every frame, in Ghoul 2 units,
  $(0, 0, -24)$ by default to match the stock animations.
  \item[Y Up to Z Up] Axis conversion, on by default.
  \item[Add as NLA Strip] Puts the animation in a new NLA strip rather than making it the active
  action. On by default.
  \item[Strip Start] First frame of the strip. $-1$ places it right after the last existing strip.
  \item[Sequence Name] Name of the action and of the \texttt{animation.cfg} sequence. Empty uses the
  file name. Several names separated by dots, e.g.
  \texttt{BOTH\_SABERFAST\_STANCE.BOTH\_STAND2.BOTH\_SABERSLOW\_STANCE}, create one action per name,
  with their strips stacked at the same frames on successive tracks.
  \item[Loop Frame] \texttt{animation.cfg} loop frame: $-1$ for no loop, 0 to loop from the start.
  \item[Start at Frame 0] Without NLA strip only: shifts the action so it starts at frame 0.
  \item[Set Scene Range \& FPS] Extends the scene's frame range to the new strips --- the
  \texttt{.gla} export samples the scene range --- and uses the file's frame rate when the strips are
  the only ones.
 \end{description}

 Strips are laid out like the \texttt{jediacademy} add-on's \texttt{.gla} import with an
 \texttt{animation.cfg}: one action per sequence (with a fake user), strips on
 ``Sequences Layer $n$'' tracks, no extrapolation, replace blending, and the sequence's loop frame,
 frame rate and length stored in the action's Ghoul 2 sequence properties when the
 \texttt{jediacademy} add-on is enabled. Only the new strips are left selected, ready to be exported.

 If another track is soloed, the new strips are not visible in the viewport until it is un-soloed;
 the \texttt{jediacademy} exports are not affected.

 \section{Export}

 Writes the armature's selected NLA strips as a dotXSI 3.0 file, in the layout of 3ds Max's
 exporter. Only the armature is exported.

 \begin{itemize}
  \item Armature: the active armature, else \texttt{skeleton\_root}, else the first armature of the
  scene.
  \item Strips: the selected strips of that armature, in timeline order. Several strips are
  concatenated end to end, and a warning says so. With no strip selected, only the skeleton (its rest
  pose, without animation) is exported, with a warning.
  \item A strip's length is the sequence length recorded by the \texttt{jediacademy} add-on (so a
  1-frame still exports 1 frame), else the length of the action range used by the strip.
  \item \texttt{face} is written as \texttt{face\_always\_}.
 \end{itemize}

 The export dialog lists the armature and the strips that will be exported. Its settings (scale,
 origin offset, axis conversion) are the inverse of the import's.

 The export reads the actions' keyframes directly: constraints, drivers, and a strip's repeat or
 scale are not taken into account.

 \section{Workflow with the Jedi Academy plugin suite}

 \begin{enumerate}
  \item Import \texttt{models/players/\_humanoid/\_humanoid.gla} with the \texttt{jediacademy}
  add-on, with ``animations: Cfg'' to get every stock sequence as NLA strips, or without animations
  to only get the skeleton.
  \item Import the \texttt{.xsi} file. Give the sequence name(s) it should have in
  \texttt{animation.cfg}.
  \item Export the \texttt{.gla} with the \texttt{jediacademy} add-on, with
  \texttt{models/players/\_humanoid/\_humanoid} as ``gla reference'' so the bone indices match the
  game's models, and the \texttt{animation.cfg} from the NLA strips.
 \end{enumerate}

 To edit a stock animation in another tool, select its strip and export it as dotXSI.

 \subsection*{Saber stances}

 For reference, the stances of the single saber styles in \texttt{animation.cfg} are
 \texttt{BOTH\_SABERFAST\_STANCE} (fast), \texttt{BOTH\_STAND2} (medium, with its random idles
 \texttt{BOTH\_STAND2IDLE1} and \texttt{BOTH\_STAND2IDLE2}) and \texttt{BOTH\_SABERSLOW\_STANCE}
 (strong). Dual and staff use \texttt{BOTH\_SABERDUAL\_STANCE} and \texttt{BOTH\_SABERSTAFF\_STANCE}.

\end{document}
